\documentclass[12pt, letterpaper]{article}
\title{Marketing Teleportation Services}
\date{\today}

\begin{document}

\maketitle
Assuming that scientists have created a teleportation device that does provably work, many people will still be cautious about using this device for various reasons. One basic claim someone may make is that such a device simply does not exist or at least work. This is easy to disprove, though, by simply showing this person someone walking into the machine in one place and immediately coming out in a different place. Another claim, which is understandable, is that it is too dangerous. This may be true, but if something does not have risks, it likely does not work at all. Furthermore, you already drive your car, fly in planes, and take medicine every day, all of which could kill you at any moment, but you are willing to take that risk. It could also be argued that if, for example, you were driving in your car for 20 minutes, that is a constant life-threatening risk that lasts for 20 minutes, whereas in a teleporter, you wont be at risk for very long. A more philosophical question would be, is that person coming out of the teleporter truly the same person that went in? If so, I do not believe many people would have a problem with it. This question relies more on whether humans have souls rather than being an issue with the machine itself. For instance, an atheist, which would ironically probably be many of the people asking this question, believe that they are purely chemicals; therefore, if the teleporter can perfectly recreate their brain, they are the same person. However, if humans do have souls, the soul must be connected to the physical body in some way, and the teleporter would be a transference of consciousness into an identical body instead of a clone, which despite having the same memories, would still have a will distinct from the original person.

\end{document}

%%% Local Variables:
%%% mode: LaTeX
%%% TeX-master: t
%%% End:
